\documentclass[12pt]{article}
\usepackage[a4paper, margin=1in]{geometry}
\usepackage{amsmath}
\usepackage{amssymb}
\usepackage{amsthm}
\usepackage{amsfonts}
\usepackage{tabularx}
\newcounter{questionnumber}
\newenvironment{question}[1][]{\stepcounter{questionnumber}\par\noindent\textbf{Question \thequestionnumber. #1}\par}{\par}
\newenvironment{solution}{\par\noindent\textbf{Solution:}\par}{\par\vspace{1em}}
\begin{document}
\section{Discussion 00a}
\begin{question}
    Convert the following English sentences into propositional logic and the following propositions into English. State whether or not each statement is true with brief justification.\\
    (a) There is a real number which is not rational.\\
    (b) All integers are natural numbers or are negative, but not both.\\
    (c) If a natural number is divisible by 6, it is divisible by 2 or it is divisible by 3.\\
    (d) $(\forall x \in \mathbb{Z}) \, (x \in \mathbb{Q})$\\
    (e) $(\forall x \in \mathbb{Z}) \, (((2 \mid x) \lor (3 \mid x)) \implies (6 \mid x))$\\
    (f) $(\forall x \in \mathbb{N}) \, ((x > 7) \implies ((\exists a, b \in \mathbb{N}) \, (a + b = x)))$
\end{question}
\begin{solution}
    \noindent
    (a) $\left( \exists x\in \mathbb{R} \right) \left( x\notin \mathbb{Q} \right) $ 
    True. $\sqrt{2}$ is a real number which is not rational.\\
    (b) $\left( \forall x\in \mathbb{Z} \right) \left( \left( \left( x\in \mathbb{N} \right) \lor \left( x<0 \right) \right) \land \lnot \left( \left( x\in \mathbb{N} \right) \land \left( x<0 \right) \right) \right) $
    True. Trivial by definition.\\
    (c) $\left( \forall x\in \mathbb{N} \right) \left( \left( 6\mid x \right) \Longrightarrow \left( \left( 2\mid x \right) \lor \left( 3\mid x \right) \right) \right) $
    True. x is divisible by 6 means it's divisible by both 2 and 3, therefore divisible by 2 or 3.\\
    (d) All integers are rational numbers. True. Any integers equals to itself devided by 1.\\
    (e) If an integer is divisible by 2 or it is divisible by 3, it is divisible by 6. False. 2 is divisible by 2 but not divisible by 6.\\
    (f) Any natural number larger than 7 is a sum of two natural numbers. True. Any natural number is itself added by 0.    
\end{solution}
\begin{question}
    Determine whether the following equivalences hold, by writing out truth tables. Clearly state whether or not each pair is equivalent.\\
    (a) $P \land (Q \lor P) \equiv P \land Q$\\
    (b) $(P \lor Q) \land R \equiv (P \land R) \lor (Q \land R)$\\
    (c) $(P \land Q) \lor R \equiv (P \lor R) \land (Q \lor R)$
\end{question}
\begin{solution}
    \noindent
    (a) It doesn't hold.\\
    \begin{tabular}{|c|c|c|c|}
        \hline
        $P$ & $Q$ & $P \land (Q \lor P)$ & $P \land Q$\\
        \hline
        T & T & T & T \\
        \hline
        T & F & T & F \\
        \hline
        F & T & F & F \\
        \hline
        F & F & F & F \\
        \hline
    \end{tabular}\\
    (b) It holds.\\
    \begin{tabular}{|c|c|c|c|c|}
        \hline
        $P$ & $Q$ & $R$ & $(P \lor Q) \land R$ & $(P \land R) \lor (Q \land R)$\\
        \hline
        T & T & T & T & T \\
        \hline
        T & T & F & F & F \\
        \hline
        T & F & T & T & T \\
        \hline
        T & F & F & F & F \\
        \hline
        F & T & T & T & T \\
        \hline
        F & T & F & F & F \\
        \hline
        F & F & T & F & F \\
        \hline
        F & F & F & F & F \\
        \hline
    \end{tabular}\\
    (c) It holds.\\
    \begin{tabular}{|c|c|c|c|c|}
        \hline
        $P$ & $Q$ & $R$ & $(P \land Q) \lor R$ & $(P \lor R) \land (Q \lor R)$\\
        \hline
        T & T & T & T & T \\
        \hline
        T & T & F & T & T \\
        \hline
        T & F & T & T & T \\
        \hline
        T & F & F & F & F \\
        \hline
        F & T & T & T & T \\
        \hline
        F & T & F & F & F \\
        \hline
        F & F & T & T & T \\
        \hline
        F & F & F & F & F \\
        \hline
    \end{tabular}
\end{solution}
\begin{question}
    Which of the following implications are always true, regardless of P? Give a counterexample for each false assertion (i.e. come up with a statement P(x,y) that would make the implication false).\\
    (a) $\forall x\forall y\operatorname{P}(x,y) \implies  \forall y\forall x\operatorname{P}(x,y)$\\
    (b) $\forall x\exists y\operatorname{P}(x,y) \implies  \exists y\forall x\operatorname{P}(x,y)$\\
    (c) $\exists x\forall y\operatorname{P}(x,y) \implies  \forall y\exists x\operatorname{P}(x,y)$
\end{question}
\begin{solution}
    \noindent
    (a) True.\\
    (b) False. Let $\mathrm{P}\left( x,y \right) =x<y$.\\
    (c) True.
\end{solution}
\section{Discussion 00b}
\begin{question}
    \noindent
    (a) Prove that if $n^{2}$ is odd, then $n$ must also be odd.\\
    (b) Prove that if $n^{2}$ is odd, then $n^{2}$ can be written in the form $8k+1$ for some integer $k$.
\end{question}
\begin{solution}
    \noindent
    (a) Assume that $n$ is even, then let $n = 2k, k \in \mathbb{Z}$.
    Therefore $n^2 = 4k^2$. So $n^2$ is even, which is contradictory to $n^2$ is odd.
    Therefore, $n$ is odd.\\
    (b) According to (a), $n$ is odd since $n^2$ is odd. 
    Therefore, assume that $n = 2m+1, m \in \mathbb{Z}$.
    Therefore, $n^2 = 4m^2+4m+1$.
    Let $k = \frac{m\left( m+1 \right)}{2}$, then $8k+1=4m^2+4m+1=n^2$.
\end{solution}
\begin{question}
    Prove that if there are $n\geq2$ people at a party, then at least 2 of them have the same number of friends at the party. Assume that friendships are always reciprocated: that is, if Alice is friends with Bob, then Bob is also friends with Alice.\\
(Hint: The Pigeonhole Principle states that if $n$ items are placed in $m$ containers, where $n>m$,at
least one container must contain more than one item. You may use this without proof.)
\end{question}
\begin{solution}
    Assume that all n people have different numbers of friends at the party,
    which means that each of them has 0 to $n-1$ friends.
    However, one have $n-1$ friends means that they have been friends with everyone else,
    which is impossible if another one has 0 friend.
    Therefore, at least 2 of them have the same number of friends.
\end{solution}
\begin{question}
    Suppose you have a rectangular array of pebbles, where each pebble is either red or blue. Suppose that for every way of choosing one pebble from each column, there exists a red pebble among the chosen ones.\\
    Prove that there must exist an all-red column.
\end{question}
\begin{solution}
    Assume that there is not an all-red column, which means that there is always at least one blue pebble in every column.
    Therefore, we can choose one blue pebble from each column.
    But it is contradictory to the assumption that for every way of choosing on pebble from each column, there exists a red pebble among the chosen ones.
    Therefore, there must exist an all-red column.
\end{solution}
\section{Homework 00}
\begin{question}
    In parts (a)-(b), convert the English sentences into propositional logic. In parts (c) - (d), convert the propositions into English. For parts (b) and (d), use the notation $a | b$ to denote the statement``a divides b'', and use the notation $P(x)$ to denote the statement ``x is a prime number''.\\
    (a) For every real number $k$, there is a unique real solution to $x^3 = k$.\\
    (b) If $p$ is a prime number, then for any two natural numbers $a$ and $b$, if $p$ doesn't divide $a$ and $p$ divides $ab$, then $p$ divides $b$.\\
    (c) $(\forall x, y \in \mathbb{R})\left[(xy = 0) \implies ((x = 0) \lor (y = 0))\right]$\\
    (d) $\neg\left(\exists y \in \mathbb{N}\right)\left[\left(\forall x \in \mathbb{N}\right)\left[(x > y) \implies \left((y | x) \lor P(x)\right)\right]\right]$
\end{question}
\begin{solution}
    \noindent
    (a) $\left( \forall k\in \mathbb{R} \right) \left[ \left( \exists x\in \mathbb{R} \right) \left( x^3=k \right) \land \left( \forall y,z\in \mathbb{R} \right) \left( \left( y^3=k \right) \land \left( z^3=k \right) \right) \Longrightarrow \left( y=z \right) \right] $\\
    (b) $\left( \forall p,a,b\in \mathbb{N} \right) \left[ \left( \mathrm{P}\left( p \right) \right) \Longrightarrow \left( \left( \lnot \left( p\mid a \right) \land \left( p\mid ab \right) \right) \Longrightarrow \left( p\mid b \right) \right) \right] $\\
    (c) If the product of two real number is 0, then one of them must be 0.\\
    (d) There is no natural number that divides every composite number greater than it.
\end{solution}
\section{Discussion 01a}
\begin{question}
    $\text{Prove that if }n\in\mathbb{N}\mathrm{~and~}x>0\text{, then }(1+x)^n\geq1+nx.$
\end{question}
\begin{solution}
    If $n=0$,then $\left( 1+x \right) ^n=1\ge 1=1+nx$.
    Assume that the proposition is true for $n=k$, 
    then for $n=k+1$, $\left( 1+x \right) ^{k+1}=\left( 1+x \right) ^k\left( 1+x \right) \ge \left( 1+kx \right) \left( 1+x \right) =1+\left( k+1 \right) x+kx^2\ge 1+\left( k+1 \right) x$.
    Therefore, the statement is true for all $n\in\mathbb{N}$.
\end{solution}
\begin{question}
    Suppose that the sequence $a_1, a_2, \ldots$ is defined by $a_1 = 1$ and $a_{n+1} = 3a_n^2$ for $n \geq 1$. We want to prove that
    $$a_n \leq 3^{(2^n)}$$
    for every positive integer $n$.\\
    (a) Suppose that we want to prove this statement using induction. Can we let our inductive hypothesis be simply $a_n \leq 3^{(2^n)}$? Attempt an induction proof with this hypothesis to show why this does not work.\\
    (b) Try to instead prove the statement $a_n \leq 3^{(2^n-1)}$ using induction.\\
    (c) Why does the hypothesis in part (b) imply the overall claim?
\end{question}
\begin{solution}
    \noindent
    (a) Assume that $a_k\le 3^{(2^k )}$, then$a_{k+1}=3a_{k}^{2}\le 3^{2^{k+1}+1}$.The proof does not work.\\
    (b) For $n=1$, $a_1=1\le 3^{2^1-1}=3$.
    Assume that $a_k\le 3^{2^k-1}, k \geq 1$, then $a_{k+1}=3a_{k}^{2}\le 3\left( 3^{2\left( 2^k-1 \right)} \right) =3^{2^{k+1}-1}$.
    Therefore the statement is true for all $n\in\mathbb{N}.$\\
    (c) For any $n\in\mathbb{N}$, $a_n\le 3^{2^n-1}\le 3^{2^n}$.
\end{solution}
\begin{question}
    Prove that every positive integer $n$ can be written in binary. In other words, prove that for any positive
    integer $n$, we can write

    $$n=c_k\cdot2^k+c_{k-1}\cdot2^{k-1}+\cdots+c_1\cdot2^1+c_0\cdot2^0,$$

    for some $k\in\mathbb{N}$ and $c_i\in\{0,1\}$ for all $i\leq k.$
\end{question}
\begin{solution}
    For $n=1$, $1=1\cdot 2^0$.
    Assume that $n \in \mathbb{N}$ can be written in binary for some $k\in\mathbb{N}$ and $c_i\in\{0,1\}$ for all $i\leq k.$
    Let $c_{k+1}=0$. Let $j$ be the least natural number that makes $c_j=0$.
    Let $d_i=\begin{cases}
	c_i&		i>j\\
	1&		i=j\\
	0&		i<j\\
    \end{cases}$ for all $i \leq k+1$.
    Let $m=d_{k+1}\cdot2^{k+1}+d_k\cdot2^k+\cdots+d_1\cdot2^1+d_0\cdot2^0$.
    Then $m-n=(d_{k+1}-c_{k+1})\cdot2^{k+1}+(d_k-c_k)\cdot2^k+\cdots+(d_1-c_1)\cdot2^1+(d_0-c_0)\cdot2^0$.
    Since $d_i-c_i=\begin{cases}
	0&		i>j\\
	1&		i=j\\
	-1&		i<j\\
    \end{cases}$, $m-n=2^j-(2^{j-1}+2^{j-2}+\cdots+2^1+2^0)=1$.
    Therefore, we have proven that $m=n+1$ and it can be written in binary. Then the statement is true for all positive integer $n$.
\end{solution}
\begin{question}
    Recall, the Fibonacci numbers, defined recursively as
    $F_{1}=1$, $F_{2}=1$, and $F_{n}=F_{n-2}+F_{n-1}$.\\
    Prove that every third Fibonacci number is even. For example, $F_{3}=2$ is even and $F_{6}=8$ is even.
\end{question}
\begin{solution}
    First, we prove that $F_{3k}$ is even and $F_{3k-1}, F_{3k-2}$ are odd for all positive integer $k$.
    For $k=1$, $F_1=1, F_2=1, F_3= 2$, therefore $F_1, F_2$ are odd and $F_3$ is even. Then, assume that $F_{3k}$ is even and $F_{3k-1}, F_{3k-2}$ are odd.
    Therefore, $F_{3k+1}=F_{3k}+F_{3k-1}$ is odd, $F_{3k+2}=F_{3k}+F_{3k+1}$ is odd, $F_{3k+3}=F_{3k+1}+F_{3k+2}$ is even. Therefore, we have proven that $F_{3k}$ is even and $F_{3k-1}, F_{3k-2}$ are odd for all positive integer $k$.
    Therefore, we have proven that every third Fibonacci number is even.
\end{solution}
\begin{question}
    Consider the set of jobs $J = \{1, 2, 3\}$ and the set of candidates $C = \{A, B, C\}$ with the following preferences.

\begin{tabular}{|c|c|}
    \hline
    Jobs & Candidates \\
    \hline
    1 & $A > B > C$ \\
    \hline
    2 & $B > A > C$ \\
    \hline
    3 & $A > B > C$ \\
    \hline
\end{tabular}
\begin{tabular}{|c|c|}
    \hline
    Candidates & Jobs \\
    \hline
    A & $2 > 1 > 3$ \\
    \hline
    B & $1 > 3 > 2$ \\
    \hline
    C & $1 > 2 > 3$\\
    \hline
\end{tabular}

Run the traditional propose-and-reject algorithm on this example. How many days does it take and what is the resulting pairing? (Show your work.)
\end{question}
\begin{solution}
    \begin{tabular}{|c|c|c|c|}
        \hline
        Days & A & B & C\\
        \hline
        1 & 1, 3 & 2 & \\
        \hline
        2 & 1 & 2, 3 & \\
        \hline
        3 & 1, 2 & 3 & \\
        \hline
        4 & 2 & 1, 3 & \\
        \hline
        5 & 2 & 1 & 3 \\
        \hline
    \end{tabular}
    
    It takes 5 days.
\end{solution}
\begin{question}
    \noindent
    Prove the following statements about the traditional propose-and-reject algorithm.\\
    (a) In any execution of the algorithm, if a candidate receives a proposal on day $i$, then she receives some proposal on every day thereafter until termination.\\
    (b) In any execution of the algorithm, if a candidate receives no proposal on day $i$, then she receives no proposal on any previous day $j$, $1 \leq j < i$.\\
    (c) In any execution of the algorithm, there is at least one candidate who only receives a single proposal. (Hint: use the parts above!)
\end{question}
\begin{solution}
    \noindent
    (a) First, she receives some proposal on day $i$. Then, assume that she receives some proposal on day $k(k\geq i)$. By definition, she has a job on the string and the job will again give her an offer on day $k+1$. Therefore, we have proven the statement.\\
    (b) Assume that she receives at least one proposal on day $j$, then according to (a), she receives some proposal on day $i$, which leads to contradiction.\\
    (c) On the day before the last day, there is at least one candidate who receives no proposal, otherwise the algorithm will terminate on that day.
    According to (b), this candidate receives no proposal on any previous day. So she only receives a proposal on the last day.
\end{solution}
\begin{question}
    \noindent
    By stable matching instance, we mean a set of jobs and candidates and their preference lists. For each of the following statements, indicate whether the statement is True or False and justify your answer with a short 2-3 line explanation:\\
    (a) There is a stable matching instance for $n$ jobs and $n$ candidates for $n > 1$, such that in a stable matching algorithm with jobs proposing, every job ends up with its least preferred candidate.\\
    (b) In a stable matching instance, if job $J$ and candidate $C$ each put each other at the top of their respective preference lists, then $J$ must be paired with $C$ in every stable pairing.\\
    (c) In a stable matching instance with at least two jobs and two candidates, if job $J$ and candidate $C$ each put each other at the bottom of their respective preference lists, then $J$ cannot be paired with $C$ in any stable pairing.\\
    (d) For every $n > 1$, there is a stable matching instance for $n$ jobs and $n$ candidates which has an unstable pairing where every unmatched job-candidate pair is a rogue couple or pairing.
\end{question}
\begin{solution}
    \noindent
    (a) False. If that was the case, then there would be at least one rogue couple since they both preferred each other to the current partner.\\
    (b) True. If not the case then $(J,C)$ would be a rouge couple.\\
    (c) False. Consider the case that $J_1$ and $J_2$ both prefer $C_1$ and $C_2$, and $C_1$ and $C_2$ both prefer $J_1$ to $J_2$. According to (b), $J_1$ must be paired with $C_1$, then $J_2$ will be paired with $C_2$.\\
    (d) True. Consider the case that every job and candidate have their current partner at the bottom of their preference lists.
\end{solution}
\begin{question}
    \noindent
    (a) True or False?\\
    \indent (i) If a candidate accidentally rejects a job she prefers on a given day, then the algorithm still always ends with a matching.\\
    \indent (ii) The Propose-and-Reject Algorithm never produces a candidate-optimal matching.\\
    \indent (iii) If the same job is last on the preference list of every candidate, the job must end up with its least preferred candidate.\\
    (b) As you've seen from lecture, the jobs-proposing Propose-and-Reject Algorithm produces an employer-optimal stable matching. Let's see if the candidate have any way of improving their standing. Suppose exactly one of the candidates has the power to arbitrarily reject one proposal, regardless of which job she has on her string (if any). Construct an example that illustrates the following: for any $n \geq 2$, there exists a stable matching instance for which using this power helps every candidate, i.e. every candidate gets a better job than she would have gotten under the jobs-proposing Propose-and-Reject Algorithm.
\end{question}
\begin{solution}
    \noindent
    (a)(i)F(ii)F(iii)F\\
    (b)Consider the case that every $J_i$'s top 2 preferences are $C_i$ and $C_{(i+1)\mod n}$,
    and every $C_i$'s top 2 preferences are $J_{(i-1)\mod n}$ and $J_i$.
    Every candidate can get a more preferred job if $C_1$ rejects $J_1$ in the first round and accepts $J_n$ eventually.
\end{solution}
\section{Homework 01}
\begin{question}
    \noindent
    $\begin{array}{l}
    \text { Decide whether each of the following logical equivalences is correct and justify your answer. } \\
    (a) \quad \forall x(P(x) \wedge Q(x)) \stackrel{?}{=} \forall x P(x) \wedge \forall x Q(x) \\
    (b) \quad \forall x(P(x) \vee Q(x)) \stackrel{?}{=} \forall x P(x) \vee \forall x Q(x) \\
    (c) \quad \exists x(P(x) \vee Q(x)) \stackrel{?}{=} \exists x P(x) \vee \exists x Q(x) \\
    (d) \quad \exists x(P(x) \wedge Q(x)) \stackrel{?}{=} \exists x P(x) \wedge \exists x Q(x)
    \end{array}$
\end{question}
\begin{solution}
    \noindent
    (a) Correct. They imply each other.\\
    (b) Not correct. For every $x\in {0,1}$, let $\mathrm{P}(x): x=1, \mathrm{Q}(x): x=0$.\\ 
    (c) Correct. They imply each other.\\
    (d) Not correct. For every $x\in {0,1}$, let $\mathrm{P}(x): x=1, \mathrm{Q}(x): x=0$.
\end{solution}
\begin{question}
    \noindent
    For each of the following, either prove the statement, or disprove by finding a counterexample.\\
    (a) $(\forall n\in \mathbb{N})$ if $n$ is odd then $n^2+4n$ is odd.\\
    (b) $(\forall a,b\in \mathbb{R})$ if $a+b\leq 15$ then $a\leq 11$ or $b\leq 4$.\\
    (c) $(\forall r\in \mathbb{R})$ if $r^2$ is irrational, then $r$ is irrational.\\
    (d) $(\forall n\in \mathbb{Z}^+)$ $5n^3>n!$. (Note: $\mathbb{Z}^+$ is the set of positive integers)\\
    (e) The product of a non-zero rational number and an irrational number is irrational.
\end{question}
\begin{solution}
    \noindent
    (a) True. $n^2+4n=n\left( n+4 \right) $ is odd.\\
    (b) True. Assume that $a>11$ and $b>4$, then $a+b>15$, which leads to contradiction.\\
    (c) True. Assume that $r$ is rational, then $r^2$ must be rational, which leads to contradiction.
    (d) False. Take $n=10$.\\
    (e) True. Assume that the product is irrational, then divide it by the rational factor, then the irrational factor must be rational, which leads to contradiction.
\end{solution}
\begin{question}
    \noindent
    (a) Let $p > 3$ be a prime. Prove that $p$ is of the form $3k + 1$ or $3k - 1$ for some integer $k$.\\
    (b) Twin primes are pairs of prime numbers $p$ and $q$ that have a difference of 2. Use part (a) to prove that 5 is the only prime number that takes part in two different twin prime pairs.
\end{question}
\begin{solution}
    \noindent
    (a) Assume that $p=3k$ for some integer k.Then $p$ has 3 as its factor, which contradicts with the fact that p is a prime.\\
    (b) Assume that there is a prime number $p$ greater than 5 that takes part in two different twin prime pairs.
    If $p=3k-1$, then $p-2=3k-3\geq 6$ has 3 as its factor. If $p=3k+1$, then $p+2=3k+3\geq 9$ has 3 as its factor.
    Both cases lead to contradiction, therefore, 5 is the only prime number that takes part in two different twin prime pairs.
\end{solution}
\begin{question}
    Suppose that there are $2n+1$ airports, where $n$ is a positive integer. The distances between any two airports are all different. For each airport, exactly one airplane departs from it and is destined for the closest airport. Prove by induction that there is an airport which has no airplanes destined for it.
\end{question}
\begin{solution}
    For $n=1$, name the 3 airports $a_1, a_2, a_3$. Without loss of generality, assume that the distance between $a_1$ and $a_2$ is the shortest.
    Therefore, planes from $a_1$ and $a_2$ exchange their position, and the plane from $a_3$ left $a_3$ as well. Therefore, $a_3$ has no airplanes destined for it.
    Assume that the statement holds true for any $n\leq k$. For $n=k+1$, remove the closest 2 airports and there are $2k+1$ airports left, so that there is at least 1 airport for which no airplanes destined, which we name $a_m$.
    Then, add the 2 airports back, and the airplanes from the 2 airports only exchange their positions.
    The closest airport for other airplanes is either the new airports or their original destination, so still no airplanes will destined for $a_m$, so $a_m$ still has no airplanes destined for it.
    Therefore, the statement holds true for all positive integer $n$.
\end{solution}
\begin{question}
    Your "friend" Stanley Ford suggests you play the following game with him. You each start with a single stack of $n$ coins. On each of your turns, you select one of your stacks of coins (that has at least two coins) and split it into two stacks, each with at least one coin. Your score for that turn is the product of the sizes of the two resulting stacks (for example, if you split a stack of 5 coins into a stack of 3 coins and a stack of 2 coins, your score would be $3\cdot2=6).$ You continue taking turns until all your stacks have only one coin in them. Stan then plays the same game with his stack of $n$ coins, and whoever ends up with the largest total score over all their turns wins.
    
    Prove that no matter how you choose to split the stacks, your total score will always be $\frac{n(n-1)}2.$ (This means that you and Stan will end up with the same score no matter what happens, so the game is rather pointless.)
\end{question}
\begin{solution}
    For $n=2$, the total score can only be $1\cdot 1=1=\frac{1\cdot 2}{2}$.
    Assume that the statement holds true for $n=k$. 
    For $n=k+1$, split the stack into $m$ coins and $k+1-m$ coins, and it gives the player $m(k+1-m)$ points.
    The $m$-coins stack can produce $\frac{m\left( m-1 \right)}{2}$ points, and the $(k+1-m)$-coins stack can produce $\frac{\left( k+1-m \right) \left( k-m \right)}{2}$ points.
    Add all the scores up, and the total score will be $m(k+1-m)+\frac{m\left( m-1 \right)}{2}+\frac{\left( k+1-m \right) \left( k-m \right)}{2}=\frac{k\left( k+1 \right)}{2}$ points.
\end{solution}
\begin{question}
    Pacman is walking on an infnite 2D grid. He starts at some location $(i,j)\in\mathbb{N}^2$ in the first quadrant,
    and is constrained to stay in the first quadrant (say, by walls along the $x$ and y axes).

    Every second he does one of the following (if possible):

    (i) Walk one step down, to $( i, j-1) .$

    (ii) Walk one step left, to $(i-1,j).$

    For example, if he is at (5,0), his only option is to walk left to (4,0); if Pacman is instead at (3,2), he could walk either to (2,2) or (3,1).
    Prove by induction that no matter how he walks, he will always reach (0,0) in finite time.\\
    (Hint: Try starting Pacman at a few small points like (2,1) and looking all the different paths he could take to reach (0,0). Do you notice a pattern in the number of steps he takes? Try to use this to strengthen the inductive hypothesis.)
\end{question}
\begin{solution}
    Strenthen the inductive hypothesis to that Pacman takes exactly $i+j$ steps to reach $(0,0)$. 
    When Pacman starts at $(i,0)$, it takes him $i$ steps to reach $(0,0)$ because he can only go down staright.
    For the same reason, it takes him exactly $j$ steps to reach $(0,0)$ when he starts at $(0,j)$.
    Assume that when Pacman starts at $(i, j-1)$ or $(i-1, j)$ $(i, j > 0)$, he can only reach $(0,0)$ in exactly $(i+j-1)$ steps.
    When he starts at $(i,j)$, whether he reaches $(i,j-1)$ or $(i-1,j)$ one step later, he can only reach $(0,0)$ in exactly $(i+j-1)$ steps.
    So he can only reach $(0,0)$ in $i+j-1+1=i+j$ steps, which is a finite period of time.
\end{solution}
\section{Homework 02}
\begin{question}
    \noindent
    Suppose that preferences in a stable matching instance are universal: all $n$ jobs share the preferences $C_1 > C_2 > \cdots > C_n$ and all candidates share the preferences $J_1 > J_2 > \cdots > J_n$.\\
    (a) What pairing do we get from running the algorithm with jobs proposing? Can you prove this happens for all $n$?\\
    (b) What pairing do we get from running the algorithm with candidates proposing?\\
    (c) What does this tell us about the number of stable pairings?
\end{question}
\begin{solution}
    \noindent
    (a) Every $C_i$ is paired with $J_i$. When $n=1$, it's trivial.
    Assume that the statement holds true when $n=k$ . 
    When $n=k+1$, in the first step, every job proposes to $C_1$ and is rejected except $J_1$, leaving $k$ jobs and $k$ candidates sharing the same preference list.
    Therefore accoreding to the assumption, every $J_i$ will be paired with $C_i$.
    Therefore, we have proven that this happens for all $n$.\\
    (b) Same as (a).\\
    (c) The job-optimal pairing and the job-pessimal pairing is the same, therefore there is only one stable pairing.
\end{solution}
\begin{question}
    Prove that for every even $n \geq 2$, there exists an instance of the stable matching problem with $n$ jobs and $n$ candidates such that the instance has at least $2^{n/2}$ distinct stable matchings.
\end{question}
\begin{solution}
    When $n=2$, let the preference of $J_1$ be $C_1, C_2$, the one of $J_2$ be $C_2, C_1$, the one of $C_1$ be $J_2, J_1$, and the one of $C_2$ be $J_1, J_2$.
    Therefore, there exists $2^{2/2}=2$ stable matching: $(J_1, C_1), (J_2, C_2)$ and $(J_1, C_2), (J_2, C_1)$.
    Assume that the statement holds true for $n=2k$, then add 2 more jobs $J_{k+1}, J_{k+2}$ and 2 more candidates $C_{k+1}, C_{k+2}$.
    Add them to the bottom of the first $2k$ jobs and candidates' preference lists, so that there will be no new rouge couples.
    Then, let the top 2 of the preference of $J_{k+1}$ be $C_{k+1}, C_{k+2}$, the one of $J_{k+2}$ be $C_{k+2}, C_{k+1}$, the one of $C_{k+1}$ be $J_{k+2}, J_{k+1}$, and the one of $C_{k+2}$ be $J_{k+1}, J_{k+2}$.
    Therefore, there will independently be $2^k$ stable matchings among the first $2k$ jobs and candidates and 2 stable matchings among the last 2 jobs and candidates.
    Multiply them up, and we can get at least $2^{k+1}$ distinct stable matchings.
\end{solution}
\begin{question}
    \noindent
    (a) In the notes, we show that the stable matching algorithm terminates in at most $n^2$ days. Prove the following stronger result: the stable matching algorithm will always terminate in at most $(n-1)^2 + 1 = n^2 - 2n + 2$ days.\\
    (b) Provide a set of preference lists for 4 jobs and 4 candidates that will result in the upper bound from part (a) when running the Propose-and-Reject algorithm. Verify this by running the Propose-and-Reject algorithm on your preference lists.
\end{question}
\begin{solution}
    \noindent
    (a) Before termination, not all candidates have an offer, which means at most $(n-1)$ candidates have the chance to reject. 
    And, every candidate can reject only $(n-1)$ offers. Therefore, there will be at most $(n-1)^2$ days when a job is rejected, which is before termination.
    And there will be an extra day for termination. Therefore, the algorithm will always terminate in at most $(n-1)^2+1$ days.\\
    (b)
    \begin{tabular}{|c|c|}
        \hline
        $J_1$ & $C_1>C_2>C_3>C_4$\\
        \hline
        $J_2$ & $C_2>C_3>C_1>C_4$\\
        \hline
        $J_3$ & $C_3>C_1>C_2>C_4$\\
        \hline
        $J_4$ & $C_1>C_2>C_3>C_4$\\
        \hline
    \end{tabular}
    \begin{tabular}{|c|c|}
        \hline
        $C_1$ & $J_2>J_3>J_4>J_1$\\
        \hline
        $C_2$ & $J_3>J_4>J_1>J_2$\\
        \hline
        $C_3$ & $J_4>J_1>J_2>J_3$\\
        \hline
        $C_4$ & $J_1>J_2>J_3>J_4$\\
        \hline
    \end{tabular}\\
    \begin{tabular}{|c|c|c|c|c|}
        \hline
        Day & $C_1$ & $C_2$ & $C_3$ & $C_4$\\
        \hline
        1 & $J_1,J_4$ & $J_2$ & $J_3$ & \\
        \hline
        2 & $J_4$ & $J_1,J_2$ & $J_3$ & \\
        \hline
        3 & $J_4$ & $J_1$ & $J_2,J_3$ & \\
        \hline
        4 & $J_3,J_4$ & $J_1$ & $J_2$ & \\
        \hline
        5 & $J_3$ & $J_1,J_4$ & $J_2$ & \\
        \hline
        6 & $J_3$ & $J_4$ & $J_1,J_2$ & \\
        \hline
        7 & $J_2,J_3$ & $J_4$ & $J_1$ & \\
        \hline
        8 & $J_2$ & $J_3,J_4$ & $J_1$ & \\
        \hline
        9 & $J_2$ & $J_3$ & $J_1,J_4$ & \\
        \hline
        10 & $J_2$ & $J_3$ & $J_4$ & $J_1$\\
        \hline
    \end{tabular}
\end{solution}
\begin{question}
    In the stable matching problem, suppose that some jobs and candidates have hard requirements and might not be able to just settle for anything. In other words, each job/candidate prefers being unmatched rather than be matched with those below a certain point in their preference list. Let the term "entity" refer to a candidate/job. A matching could ultimately have to be partial, i.e., some entities would and should remain unmatched.

    Consequently, the notion of stability here should be adjusted a little bit to capture the autonomy of both jobs to unilaterally fire employees and/or employees to just walk away. A matching is stable if
    \begin{itemize}
    \item there is no matched entity who prefers being unmatched over being with their current partner;
    \item there is no matched/filled job and unmatched candidate that would both prefer to be matched with each other over their current status;
    \item there is no matched job and matched candidate that would both prefer to be matched with each other over their current partners; and
    \item similarly, there is no unmatched job and matched candidate that would both prefer to be matched with each other over their current status;
    \item there is no unmatched job and unmatched candidate that would both prefer to be with each other over being unmatched.
    \end{itemize}
    (a) Prove that a stable pairing still exists in the case where we allow unmatched entities.

    (HINT: You can approach this by introducing imaginary/virtual entities that jobs/candidates ``match'' if they are unmatched. How should you adjust the preference lists of jobs/candidates, including those of the newly introduced imaginary ones for this to work?)\\
    (b) As you saw in the lecture, we may have different stable matchings. But interestingly, if an entity remains unmatched in one stable matching, they must remain unmatched in any other stable matching as well. Prove this fact by contradiction.
\end{question}
\begin{solution}
    \noindent
    (a) Introduce one robot for each candidate and each job, and have all robots set their owner at the top of their preference lists, and have the order of others randomly.
    Replace the ``unmatched'' status in each real entity's preference lists with their own robot, and have other robots at the bottom of their preference lists. 
    Then, find a stable pairing by the propose-and-reject algorithm.
    After that, delete all robots and leave these who have paired with their robots unmatched.
    By definition, the original stable pairing doesn't include rouge couples.
    Therefore, no rouge couples will be in the new pairing. 
    No matched entity will prefer being unmatched,
    because if that were to happen, that would mean the entity prefer their robot than their current partner, therefore produce a rouge couple in the original stable pairing, which leads to contradiction by definition.
    Therefore, we have proven that the new pairing is stable.\\
    (b) Without the loss of generality, we assume that a job $J_1$ is unmatched in one stable matching and is paired with $C_1$ in another.
    Therefore, when $J_1$ is unmatched, $C_1$ is paired with another job $J_2$, since it prefers $J_1$ to being unmatched.
    When $J_1$ paired with $C_1$, $J_2$ must be paired with another candidate $C_2$, since it prefers $C_1$ to being unmatched.
    The process continues. And by induction, it implies that there are infinite jobs, which leads to contradiction.
    Therefore, we have proven that an entity cannot remains unmatched in one stable matching and matched in another.
\end{solution}
\end{document}